\section{The Arrow} \label{sec:arrow}

To make distinctions is to count them. The counting is discrete, one then two then three, and the structure of the
count is the integers. This is the base, and it is discrete throughout. The continuous real line appears only later as
a smoothed-out, large-scale approximation. So when we ask what gives time its direction, the answer must be a discrete
order, not the order of the continuum.

\begin{principle}[The arrow is the order of the integers]
Among the discrete normed rings the world is built from, only the integers $\mathbb{Z}$ admit a total order. The
others, the Gaussian integers and their successors, all contain an element $i$ with $i^2 = -1$, which no ordered ring
allows. The one-way order of the integers is the arrow of time.
\end{principle}

The argument is a short theorem. In any ordered ring, every square is non-negative, because a positive times a
positive is positive and a negative times a negative is positive. The integers satisfy this and carry the familiar
order $\cdots < -1 < 0 < 1 < \cdots$. But the moment we build outward, at the very first doubling (Appendix~\ref{app:tower}),
a new element $i$ appears with $i^2 = -1$. Since $-1 < 0$, no order can place $i$, and the structure becomes
unorderable. The same holds at every higher rung. So of all the structures in the construction, the integers are the
unique one with a genuine before and after.

An order is one-way. The relation $a < b$ is not the relation $b < a$. So the integer count carries an intrinsic
direction, and one cannot count both ways at once. That one-wayness is the arrow. It is not an extra ingredient added
to the theory. It is the order of the counting itself, and it survives precisely because the integers, the count,
remain orderable while everything they build does not.

There is also a floor and no ceiling. The vacuum, the $0$ of the tone, is the least, the absence, with nothing beneath
it, while the count climbs without bound. So the available direction is one-way, upward and outward from the floor.
This is the shape of time. The past is small and few, the seed. The future is large and open, the growing edge. One
cannot return below the floor, so distinction can only accumulate.

A clean reading follows. Time is the ordered part, the integer count. The directions that are unordered, the ones
carrying $i^2 = -1$, are the space-like directions, symmetric, with no intrinsic before and after between their
points. This is why time and space feel different, and the difference is exact rather than stipulated, time is the
ordered integer order, the space-like directions are the unordered ones. We will see in Section~\ref{sec:spacetime}
that the dimension of physical space comes from a separate fact, the rank of $D_4$, and is not to be confused with the
seven unordered octonion units, which are an internal structure. The point here is only the existence and the
one-wayness of time, and that is the integer order.

A remark for later. The reversible local law we build in Section~\ref{sec:law} has no local arrow at all. It is
exactly $CPT$ symmetric, and indeed exactly time-reversal symmetric in its bulk. The global arrow, the actual
forward of the world, is the growth of the substrate from its seed, the integer order made dynamical. Reversibility in
the small and an arrow in the large are consistent, and both are present here.
